\documentclass[11pt,a4paper]{article}
\usepackage[T1]{fontenc}
\usepackage{lmodern,amsmath,amssymb,amsthm,mathtools,booktabs,microtype}
\usepackage[margin=25mm]{geometry}
\usepackage[colorlinks=true,linkcolor=blue,urlcolor=blue,citecolor=blue]{hyperref}
\setlength{\emergencystretch}{3em}
\newtheorem{theorem}{Theorem}[section]
\newtheorem{lemma}[theorem]{Lemma}
\newcommand{\Q}{\mathbb Q}\newcommand{\Z}{\mathbb Z}
\title{A 196-variable quartic with 176 quadratic summands}
\author{}\date{15 November 2026}
\begin{document}\maketitle
\begin{abstract}
Complete formal-parameter ratio equations imply three reciprocal elliptic
curve equations already present in a defining system. Removing their six
rational components gives a quartic in 196 total variables with 176
primitive quadratic summands and 2853 monomials. The rational zero set is
unchanged for every fixed coefficient parameter. The degree, exact
coefficient height, arithmetic premises and sufficient defining threshold
are unchanged. The count of variables is not reduced.
\end{abstract}

\section{The polynomial specification}
Let $K=\Q(\rho)$, $\rho^2=2$, $D=11-6\rho$, and let
$E:y^2=x^3-Dx$. Use the 196-coordinate system of~\cite{Laurent},
with input $X$ and fixed integer parameter $N\ge1$. Its unbalanced
coefficient parameter is $A=149^{2N+1}$; its balanced parameter is
$C=149^N$. Delete exactly the rows
\[
\begin{gathered}
\texttt{pair1.curve.0},\quad\texttt{pair1.curve.1},\\
\texttt{pair3.curve.0},\quad\texttt{pair3.curve.1},\\
\texttt{pair4.curve.0},\quad\texttt{pair4.curve.1}.
\end{gathered}
\]
Keep every coordinate and every other primitive integral row unchanged.
The 176 retained rows are denoted $q_{N,1},\ldots,q_{N,176}$, and set
\[
                   P_N=\sum_{j=1}^{176}q_{N,j}^{\,2}.
\]
The directory
\begin{center}\small
\nolinkurl{computation/quartic/paired_redundancy196_v40}
\end{center}
contains the entire quadratic and quartic tables in both presentations,
their source bindings, producer, independent checker and certificate.
The fixed parameter is a coefficient, not a quantified variable.

\begin{theorem}\label{thm:main}
For every $N\ge1$, the rational zero set of the new quadratic system is
exactly that of the 196-coordinate parent system. Its balanced sum of
squares $P_N\in\Z[X,Y_1,\ldots,Y_{195}]$ has total degree exactly four,
integer content one, exactly 2853 coordinate monomials, and height
\[
 H(P_N)=
 11903780636513165769654680680436132271178681610029854792
 \,149^{2N}.
\]
Under precisely the arithmetic premises of the parent construction,
every sufficiently large fixed $N$ therefore makes the rational input
projection of $P_N=0$ equal to $\Z$. The same sufficient threshold works;
no numerical value of that ineffective threshold is asserted.
\end{theorem}

The parent derives from the 198-variable construction supplied in the
Conant--ChatGPT research archive~\cite{Archive}. Its original attribution
and retained arithmetic hypotheses are unchanged. The present algebraic
equivalence adds no prime, pole, modular-height or other arithmetic premise.

\section{Recovery of the removed equations}
We restate the algebraic implication from~\cite{Paired}. Suppose
$f,h,g,k,s,F,G,L,M,S,T$ belong to a field of characteristic zero,
$D\ne0$, and $\kappa,\lambda\ne0$. Assume
\begin{align*}
 F=f^2,\qquad G=g^2,\qquad
 s(f-g)&=h+k,\\
 s(h-k)&=1-D(F+fg+G),\\
 hL+\kappa fS&=0,\qquad (1-DF)L+\kappa hS=0,\\
 kM+\lambda gT&=0,\qquad (1-DG)M+\lambda kT=0.
\end{align*}
Put $U=h^2-f+Df^3$ and $V=k^2-g+Dg^3$.
The secant and tangent equations give $U=V$, by subtraction.
The primary ratio equations give $\kappa S U=0$:
multiply their second row by $h$ and subtract $1-DF$ times their
first row. The alternate ratio equations similarly give
$\lambda T V=0$. Consequently, if $(S,T)\ne(0,0)$, both
curve defects vanish. This reasoning includes $f=0$, $g=0$, tangent
pairs and zero values of either single reference ratio; it never divides
by a numerator coordinate or ratio.

In the parent system the complete reference chart~\cite{Reference}
gives finite points $H,H-P\in E(K)$ with nonzero abscissas and
\[
 S=y(H)/x(H),\qquad T=y(H-P)/x(H-P).
\]
Its fixed point is
\[
 P=\left(\frac{57-28\rho}{8},\frac{-660+417\rho}{32}\right),
\]
whose ordinate is nonzero. If $S=T=0$, both $H$ and $H-P$ would
be points of order two, so their difference $P$ would have order
dividing two, a contradiction. Thus $(S,T)\ne(0,0)$ follows from
unchanged reference equations alone.

The three retained paired-point blocks have exactly the form above.
Their nonzero scale factors are $1,2,4$, with the affine fourth-block
shift absorbed into $L,M$. Their physical reciprocal-square coordinates
are the same as in the parent. The invertible monomial coordinate changes
of~\cite{Laurent} do not alter the implication. Its base units $R,w_0,w_1$
remain nonzero on the reduced system: the unchanged scalar product and
forced-prime rows give $R>0$, and the unchanged rows $q_jw_j=1$ give
$w_j\ne0$. Therefore the same inverse coordinate changes apply before
using any of the three removed curve equations.

Each omitted field curve equation, hence each of its two rational
components, now follows from the retained rows. The reverse implication
is immediate because rows were only deleted. This proves equality of
the whole rational zero sets, with every original exceptional stratum
included. It also proves that the sufficient defining parameter is
unchanged.

\section{Degree, height and exact verification}
The input $X$ occurs only in the primitive row defining $X^2$.
Its square contributes coefficient one to $X^4$, and no other retained
row contains $X$. Hence $P_N$ has exact degree four and integer content
one for every fixed parameter.

The coefficient assertions are supplied with finite exact certificates.
The checker expands each quadratic square by a symmetric-pair algorithm,
independently of the producer's ordered-pair multiplication. It authenticates
both complete parent tables, checks every retained row by name and source
index, and reconstructs the six removed physical curve components using
$D=11-6\rho$. It also checks the two universal curve-defect identities
coefficientwise over the polynomial ring.

For the balanced family each coordinate monomial has coefficient
$p(C)=a_0+a_1C+a_2C^2$. There are 2853 such monomials. If $a_d$ is its
leading nonzero coefficient, the certificate checks
\[
 |a_d|149^d>\sum_{j<d}|a_j|149^j.
\]
Thus no monomial disappears at any $C\ge149$. With
\[
 K_0=11903780636513165769654680680436132271178681610029854792,
\]
the checker also verifies, for every coefficient profile,
\[
 \frac{|a_0|}{149^2}+\frac{|a_1|}{149}+|a_2|\le K_0.
\]
This implies $|p(C)|\le K_0 C^2$ for all $C\ge149$.
The monomial with coordinate indices $(15,16,152,152)$ has coefficient
exactly $K_0C^2$, proving attainment and the asserted height.
These coefficient checks certify finite algebraic statements; they do
not certify the independent arithmetic premises of the defining theorem.

\begin{thebibliography}{9}
\bibitem{Laurent}\emph{Quartics in 196 and 258 variables by elimination
after monomial rescaling}, project manuscript, 15 November 2026,
\nolinkurl{papers/quartic/laurent_elimination_20261115_v38.tex}.
\bibitem{Paired}\emph{Recovery of elliptic curve equations from complete ratio equations}, project manuscript, 15 November 2026,
\nolinkurl{papers/elliptic_charts/paired_ratio_redundancy_20261115_v40.tex}.
\bibitem{Reference}\emph{A complete five-coordinate elliptic reference
chart}, corrected project manuscript, 14 November 2026,
\nolinkurl{papers/elliptic_charts/compressed_reference_20261114_v27.tex}.
\bibitem{Archive}\emph{A quartic in 198 variables}, supplied research
archive, ChatGPT in collaboration with Gabriel Conant's research request,
14 November 2026,
\nolinkurl{sources/h10_quartic198/arithmetic/inverse_vieta_198.tex}.
\end{thebibliography}
\end{document}
